\documentclass[12pt]{article}
\usepackage{bmc}
\newcommand{\bandname}{Grades 4--5}
\newcommand{\packetid}{NN-45-v1}
\newcommand{\barlabel}[3]{\begin{tabular}{@{}c@{}}\chocbar{#1}{#2}{#3}\\[2mm] #1 by #2\end{tabular}}
\newcommand{\countrec}{\begin{tabular}{@{}l@{}}dots\ \rule{10mm}{0.4pt}\\[2.5mm] lines\ \rule{10mm}{0.4pt}\\[2.5mm] regions\ \rule{10mm}{0.4pt}\end{tabular}}
\begin{document}
\fontsize{12}{16}\selectfont

% ---------------- page 1 ----------------
A break splits one piece into two pieces along a straight line between squares. Break one piece at a time, and never stack pieces. In the game ``last break wins,'' two players take turns making one break, and whoever makes the last break wins.
\gap[3mm]
\begin{center}\breakpicture[6]\end{center}
\gap[5mm]

\prob Play ``last break wins'' on the 3 by 3 bar and the 2 by 4 bar. Play each bar twice, switching who goes first, and circle whether the 1st or 2nd player won. Can the losing player ever win by breaking in a different way?
\gap[4mm]
\begin{center}
\begin{tabular}{@{}m{38mm}@{\hspace{6mm}}m{52mm}@{\hspace{6mm}}m{52mm}@{}}
\barlabel{3}{3}{7} & Game 1:\enspace 1st\qquad 2nd & Game 2:\enspace 1st\qquad 2nd\\
\noalign{\vspace{6mm}}
\barlabel{2}{4}{7} & Game 1:\enspace 1st\qquad 2nd & Game 2:\enspace 1st\qquad 2nd
\end{tabular}
\end{center}
\writelines{3}
\gap[8mm]

\prob What are the fewest breaks and the most breaks that can split the 3 by 4 bar into single squares? Explain why.
\gap[3mm]
\begin{tabular}{@{}m{40mm}@{\hspace{8mm}}m{60mm}@{}}
\barlabel{3}{4}{7} & fewest\ \rule{14mm}{0.4pt}\qquad most\ \rule{14mm}{0.4pt}
\end{tabular}
\writelines{5}

\newpage
% ---------------- page 2 ----------------
\prob Decide who wins ``last break wins'' on a 7 by 9 bar, an 8 by 13 bar and a 25 by 40 bar. Find a rule that tells who wins on any bar, and explain why it works. Who wins if the game starts with a 2 by 3 bar and a 2 by 4 bar on the table, and each move breaks any one piece?
\gap[6mm]
\begin{center}
\renewcommand{\arraystretch}{1.6}
\begin{tabular}{@{}l@{\hspace{10mm}}l@{\hspace{24mm}}l@{\hspace{10mm}}l@{}}
7 by 9 & 1st\qquad 2nd & 25 by 40 & 1st\qquad 2nd\\
8 by 13 & 1st\qquad 2nd & 2 by 3 and 2 by 4 & 1st\qquad 2nd
\end{tabular}
\end{center}
\gap[2mm]
\writelines{12}

\newpage
% ---------------- page 3 ----------------
In 2-Sprigs, each dot starts with two tokens beside it. Players take turns drawing a line from a dot to another dot or back to the same dot. A line may bend, but it may not cross itself or another line, or touch any other dot. Each end of a line uses up one token at its dot, so a loop needs a dot with two tokens. After each line, the other player takes away the tokens it used. Whoever draws the last line wins.
\gap[3mm]
\begin{center}\sprigsrules\end{center}
\gap[4mm]

\prob Play 2-Sprigs on these boards. Write how many moves each game lasted, and circle who won.
\gap[3mm]
\begin{center}
\recboard{\board{88}{66}{\Lthree}}{\movesrec\qquad\winrec}\hfill
\recboard{\board{88}{66}{\Lthree}}{\movesrec\qquad\winrec}
\\[4mm]
\recboard{\board{88}{72}{\Lfour}}{\movesrec\qquad\winrec}\hfill
\recboard{\board{88}{72}{\Lfour}}{\movesrec\qquad\winrec}
\end{center}

\newpage
% ---------------- page 4 ----------------
\prob Can a game of 2-Sprigs end while some dot still has a token? Try to make it happen on these boards. Show how, or explain why it can't happen. Then say who wins 2-Sprigs with any number of dots.
\gap[4mm]
\begin{center}
\board{90}{90}{\Lfive}\hspace{14mm}\board{80}{106}{\Lsixtall}
\end{center}
\gap[2mm]
\writelines{9}

\newpage
% ---------------- page 5 ----------------
In Sprouts, each spot starts with three tokens. Players take turns drawing a line from a spot to another spot or back to the same spot. A line may bend, but it may not cross itself or another line, or touch any other spot. Each end of a line uses up one token at its spot, so a loop needs a spot with two tokens. After each line, the other player takes away the tokens it used, and the player who drew it adds a new spot on the line and gives it one token. Whoever makes the last move wins.
\gap[3mm]
\begin{center}\sproutsrules\end{center}
\gap[4mm]

\prob Work together to make a Sprouts game that starts with 2 spots last as many moves as you can, and another end after as few moves as you can. Write how many moves each game lasted. Can the players' choices change who wins Sprouts?
\gap[3mm]
\begin{center}
\recboard{\board{91}{96}{\Ltwospot}}{\movesrec}\hfill
\recboard{\board{91}{96}{\Ltwospot}}{\movesrec}
\end{center}
\writelines{2}

\newpage
% ---------------- page 6 ----------------
\prob Work together again, starting with 3 spots: make one game last as many moves as you can, and another end after as few moves as you can. Write how many moves each game lasted.
\gap[4mm]
\begin{center}
\recboard{\board{186}{100}{\Lthreespotbig}}{\movesrec}
\\[4mm]
\recboard{\board{186}{100}{\Lthreespotbig}}{\movesrec}
\end{center}

\newpage
% ---------------- page 7 ----------------
\prob What is the largest number of moves a game of Sprouts that starts with 10 spots could last? Explain why it can never last longer.
\writelines{5}
\gap[8mm]

Brussels Sprouts starts with crosses. An arm of a cross is free when no line is attached to it. Players take turns drawing a line that joins two free arms, which may belong to the same cross. The line may not cross itself or another line, or touch any other cross. The player then draws a short bar across the new line, making a new cross with two new free arms. Whoever makes the last move wins.
\gap[3mm]
\begin{center}\brusselsrules\end{center}
\gap[4mm]

\prob Play Brussels Sprouts on these boards, and write how many moves each game lasted. Then play more games with 2 crosses on a whiteboard, trying to make one end sooner or later than the others.
\gap[3mm]
\begin{center}
\recboard{\crossboard{70}{72}{\Lone}}{\movesrec}\hfill
\recboard{\crossboard{112}{72}{\Ltwo}}{\movesrec}
\end{center}

\newpage
% ---------------- page 8 ----------------
\prob How many moves will a game of Brussels Sprouts that starts with 3 crosses last? Decide before you play, then play on this board to check. How many moves would a game that starts with 10 crosses last?
\gap[4mm]
\begin{center}
\crossboard{180}{136}{{(-38,-28),(38,-28),(0,34)}}
\end{center}
\gap[2mm]
\writelines{5}

\newpage
% ---------------- page 9 ----------------
\prob Count the dots, the lines and the regions in each drawing below. Count each piece of line from one dot to the next as one line, and count the outside as a region. When a drawing has crosses, count each cross as a dot and do not count the free arms. Find a rule that connects the three numbers, and check it on your finished games from Problem 9. Then use your rule to explain why every game of Brussels Sprouts that starts with 2 crosses lasts the same number of moves.
\gap[6mm]
\begin{center}
\begin{tabular}{@{}c@{\hspace{6mm}}c@{\hspace{6mm}}c@{\hspace{6mm}}c@{}}
\begin{tikzpicture}[mm,scale=1.35,baseline=(current bounding box.center)]
  \draw[play] (0,-7) -- (12,0) -- (24,0); \draw[play] (12,0) -- (12,13); \draw[play] (24,0) -- (24,-12);
  \foreach \p in {(0,-7),(12,0),(24,0),(12,13),(24,-12)} \fill \p circle (1.5);
\end{tikzpicture}
&
\begin{tikzpicture}[mm,scale=1.35,baseline=(current bounding box.center)]
  \draw[play] (0,0) rectangle (20,20); \draw[play] (0,0) -- (20,20);
  \foreach \p in {(0,0),(20,0),(20,20),(0,20)} \fill \p circle (1.5);
\end{tikzpicture}
&
\begin{tikzpicture}[mm,scale=1.35,baseline=(current bounding box.center)]
  \draw[play] (0,0) .. controls (6,9) and (14,9) .. (20,0);
  \draw[play] (0,0) .. controls (6,-9) and (14,-9) .. (20,0);
  \draw[play] (20,0) .. controls (33,12) and (33,-12) .. (20,0);
  \foreach \p in {(0,0),(20,0)} \fill \p circle (1.5);
\end{tikzpicture}
&
\begin{tikzpicture}[mm,baseline=(current bounding box.center)]
  % a finished game that started with 1 cross (3 moves); checked numerically
  \cross[4.5]{(0,0)}
  \draw[play] (0,4.5) .. controls (0,21) and (21,0) .. (4.5,0);
  \draw[line width=1.3pt] (5.26,5.26) -- (11.62,11.62);
  \draw[play] (0,-4.5) .. controls (0,-21) and (-21,0) .. (-4.5,0);
  \draw[line width=1.3pt] (-5.26,-5.26) -- (-11.62,-11.62);
  \draw[play] (11.62,11.62) .. controls (16.6,16.6) and (27,11) .. (27,0)
    .. controls (27,-17) and (17,-27) .. (0,-27)
    .. controls (-13,-27) and (-17,-17) .. (-11.62,-11.62);
  \draw[line width=1.3pt] (22.5,0) -- (31.5,0);
\end{tikzpicture}
\\[4mm]
\countrec & \countrec & \countrec & \countrec
\end{tabular}
\end{center}
\gap[4mm]
\writelines{9}

\end{document}
